\documentclass[12pt]{article}
\usepackage[ukrainian,shorthands=off]{babel}   % shorthands=off: символ " не активний (потрібно для міток "..." у TikZ всередині \zadtask)
\usepackage{fontspec}
% --- НАЛАШТУВАННЯ ШРИФТІВ ---
% Шрифти Schoolbook-*.otf лежать у корені проєкту. Шлях підбирається автоматично:
% ../ (компіляція з папки теми), ./ (корінь проєкту), ../../ (підпапки теми 28); інакше -- TeX Gyre Schola.
\newcommand{\nmtSchoolbook}[1]{\setmainfont{Schoolbook-Regular.otf}[
    Path = #1,
    BoldFont = Schoolbook-Bold.otf,
    ItalicFont = Schoolbook-Italic.otf,
    BoldItalicFont = Schoolbook-BoldItalic.otf
]}
\IfFileExists{../Schoolbook-Regular.otf}{\nmtSchoolbook{../}}{%
 \IfFileExists{./Schoolbook-Regular.otf}{\nmtSchoolbook{./}}{%
  \IfFileExists{../../Schoolbook-Regular.otf}{\nmtSchoolbook{../../}}{%
   \setmainfont{texgyreschola}[Extension=.otf, UprightFont=*-regular, BoldFont=*-bold, ItalicFont=*-italic, BoldItalicFont=*-bolditalic]}}}
\usepackage[italic]{mathastext}
\usepackage[a4paper,margin=1.8cm,bottom=2cm,top=2cm]{geometry}
\usepackage{amsmath,amssymb}
\usepackage{tikz}
\usepackage{pgfplots}
\pgfplotsset{compat=1.18}
\usepgfplotslibrary{fillbetween}
\usetikzlibrary{calc,patterns,angles,quotes,intersections,babel,3d,shapes.symbols,shapes.geometric,decorations.pathreplacing,shadings,arrows.meta,hobby}
\usepackage{xcolor,array,fancyhdr,enumitem,multicol}
\usepackage{colortbl,multirow,diagbox}
\usepackage{tcolorbox}
\tcbuselibrary{skins,breakable}
\usepackage[hidelinks]{hyperref}
\usepackage[firstpage=false, text=@pvtr2525, scale=0.6, color=gray!12]{draftwatermark}

% Заборона розриву інлайн-формул
\binoppenalty=10000
\relpenalty=10000

\definecolor{mainGreen}{RGB}{34, 120, 64}
\definecolor{yearOrange}{RGB}{220, 100, 30}
\definecolor{yearcolor}{RGB}{220, 100, 30}
\definecolor{titleBg}{RGB}{225, 240, 220}
\definecolor{instrBg}{RGB}{255, 240, 220}
\definecolor{instrBorder}{RGB}{220, 150, 80}
\definecolor{headerblue}{RGB}{0, 102, 204}   % використовується в рисунках старої бази

\pagestyle{fancy}
\fancyhf{}
\renewcommand{\headrulewidth}{0.4pt}
\renewcommand{\footrulewidth}{0.4pt}
\fancyhead[C]{\small\color{gray!80} База завдань НМТ 2023--2026 \textendash{} Тема 28a}
\fancyhead[R]{\small\color{mainGreen}\textbf{\href{https://t.me/pvtr2525}{@pvtr2525}}}
\fancyfoot[L]{\small\color{gray!80}\href{https://t.me/pvtr2525}{ТГ: @pvtr2525}}
\fancyfoot[C]{\small\color{gray!80}\thepage}
\fancyfoot[R]{\small\color{gray!80}НМТ 2023--2026}

% --- ТАБЛИЦІ ВІДПОВІДЕЙ ---
\newcommand{\answerTable}[5]{%
\par\vspace{0.15cm}
\noindent
\begin{tabular}{|*{5}{>{\centering\arraybackslash}m{3cm}|}}
\hline
\rule[-0.15cm]{0pt}{0.6cm}\textbf{А} & \rule[-0.15cm]{0pt}{0.6cm}\textbf{Б} & \rule[-0.15cm]{0pt}{0.6cm}\textbf{В} & \rule[-0.15cm]{0pt}{0.6cm}\textbf{Г} & \rule[-0.15cm]{0pt}{0.6cm}\textbf{Д} \\
\hline
\rule[-0.2cm]{0pt}{0.75cm}#1 & \rule[-0.2cm]{0pt}{0.75cm}#2 & \rule[-0.2cm]{0pt}{0.75cm}#3 & \rule[-0.2cm]{0pt}{0.75cm}#4 & \rule[-0.2cm]{0pt}{0.75cm}#5 \\
\hline
\end{tabular}
\par\vspace{0.25cm}
}
\newcommand{\answerTableTall}[5]{%
\par\vspace{0.15cm}
\noindent
\begin{tabular}{|*{5}{>{\centering\arraybackslash}m{3cm}|}}
\hline
\rule[-0.2cm]{0pt}{0.7cm}\textbf{А} & \rule[-0.2cm]{0pt}{0.7cm}\textbf{Б} & \rule[-0.2cm]{0pt}{0.7cm}\textbf{В} & \rule[-0.2cm]{0pt}{0.7cm}\textbf{Г} & \rule[-0.2cm]{0pt}{0.7cm}\textbf{Д} \\
\hline
\rule[-0.45cm]{0pt}{1.2cm}#1 & \rule[-0.45cm]{0pt}{1.2cm}#2 & \rule[-0.45cm]{0pt}{1.2cm}#3 & \rule[-0.45cm]{0pt}{1.2cm}#4 & \rule[-0.45cm]{0pt}{1.2cm}#5 \\
\hline
\end{tabular}
\par\vspace{0.25cm}
}
\newcommand{\instructionBox}[1]{%
\par\vspace{0.3cm}
\begin{tcolorbox}[colback=instrBg, colframe=instrBorder, boxrule=0.6pt, arc=2pt,
    left=8pt, right=8pt, top=4pt, bottom=4pt]
\centering\small\bfseries #1
\end{tcolorbox}
\par\vspace{0.15cm}
}
\newcommand{\sectionTitle}[1]{%
\par\vspace{0.4cm}
\begin{tcolorbox}[colback=titleBg, colframe=titleBg, boxrule=0pt, arc=8pt,
    halign=center, fontupper=\Large\bfseries\color{mainGreen}]
#1
\end{tcolorbox}
\par\vspace{0.3cm}
}
\newcommand{\task}[2]{\noindent\textbf{#1.}\ \ #2\par\vspace{0.2cm}}
% --- НМТ-СТИЛЬ ПОЛЕ ВІДПОВІДІ ---
\newcommand{\nmtAnswerBox}{%
\par\vspace{0.2cm}\noindent
Відповідь:\ \
\begin{tabular}{@{}*{4}{|>{\centering\arraybackslash}p{0.8cm}}|@{\hspace{0.1cm}\textbf{,}\hspace{0.1cm}}*{3}{|>{\centering\arraybackslash}p{0.8cm}}|@{}}
\hline
\rule[-0.2cm]{0pt}{0.7cm} & & & & & & \\
\hline
\end{tabular}
\par\vspace{0.3cm}
}
% --- СІТКА ДЛЯ МАТЧИНГУ ---
\newcommand{\matchingGrid}{%
    \begingroup
    \renewcommand{\arraystretch}{1.15}
    \setlength{\tabcolsep}{4pt}
    \footnotesize
    \begin{tabular}{r|c|c|c|c|c|}
         \multicolumn{1}{c}{} & \multicolumn{1}{c}{\textbf{А}} & \multicolumn{1}{c}{\textbf{Б}} & \multicolumn{1}{c}{\textbf{В}} & \multicolumn{1}{c}{\textbf{Г}} & \multicolumn{1}{c}{\textbf{Д}} \\ \cline{2-6}
         \textbf{1} & \phantom{X} & \phantom{X} & \phantom{X} & \phantom{X} & \phantom{X} \\ \cline{2-6}
         \textbf{2} & & & & & \\ \cline{2-6}
         \textbf{3} & & & & & \\ \cline{2-6}
    \end{tabular}
    \endgroup
}
% --- ДОДАТКОВІ МАКРОСИ ЗБІРНИКА ---
\newcommand{\taskBlock}[1]{%
\par\vspace{0.45cm}
\noindent{\color{mainGreen}\rule{\linewidth}{0.8pt}}\par\vspace{0.12cm}
\noindent{\small\bfseries\color{mainGreen}#1}\par\vspace{0.25cm}
}
\newcounter{zad}
\newcommand{\zadtask}[1]{\stepcounter{zad}\noindent\textbf{\thezad.}\ \ #1\par\vspace{0.2cm}}
\newcommand{\zadnum}{\stepcounter{zad}\textbf{\thezad.}\ \ }
\newcounter{ans}
\newcommand{\ansitem}[1]{\stepcounter{ans}\noindent\textbf{\theans.}\ #1\par\vspace{0.07cm}}
\newcommand{\nmtyear}[1]{\hfill{\small\color{yearOrange}(НМТ #1)}}
% --- МАКРОС КОРОТКОГО РОЗВ'ЯЗКУ ---
\newcommand{\solution}[1]{%
\par\vspace{0.12cm}
\begin{tcolorbox}[colback=titleBg!45, colframe=mainGreen!55, boxrule=0.4pt, arc=2pt,
    left=6pt, right=6pt, top=3pt, bottom=3pt, breakable]
\footnotesize\textbf{Короткий розв'язок.}\ #1
\end{tcolorbox}
\par\vspace{0.12cm}
}
% Розв'язки й відповіді (є до завдань НМТ-2026): \showsolutionstrue -- показати, \showsolutionsfalse -- сховати
\newif\ifshowsolutions
\showsolutionsfalse
% --- ЗАГОЛОВКИ З ПУНКТАМИ КЛІКАБЕЛЬНОГО ЗМІСТУ ---
\setcounter{tocdepth}{2}
\newcommand{\chapterTitle}[1]{%
\clearpage\phantomsection\addcontentsline{toc}{section}{#1}%
\sectionTitle{#1}}
\newcommand{\typeTitle}[1]{%
\phantomsection\addcontentsline{toc}{subsection}{\quad #1}%
\par\vspace{0.3cm}
\noindent{\large\bfseries\color{yearOrange}#1}\par\vspace{0.08cm}
\noindent{\color{yearOrange}\rule{\linewidth}{0.4pt}}\par\vspace{0.2cm}}
\newcommand{\ansTheme}[1]{\par\vspace{0.25cm}\noindent{\bfseries\color{mainGreen}#1}\par\vspace{0.12cm}}
\newcommand{\ansType}[1]{\par\vspace{0.1cm}\noindent{\itshape\color{yearOrange}#1}\par\vspace{0.08cm}}

\begin{document}
\begingroup\setcounter{zad}{0}
% --- макроси, перенесені зі старого файлу теми ---
\newcommand{\matchingLayout}[3]{
    \noindent
    \begin{minipage}[t]{0.40\textwidth}
       
        #1
    \end{minipage}%
    \hfill
    \begin{minipage}[t]{0.28\textwidth}
        
        #2
    \end{minipage}%
    \hfill
    \begin{minipage}[t]{0.30\textwidth}
        \vspace{0pt} % Хаки для вирівнювання minipage по верху
        \begin{flushright}
        #3
        \end{flushright}
    \end{minipage}
}
\newcommand{\answerTableSmall}[5]{
\begin{tabular}{|*{5}{>{\centering\arraybackslash}m{1.65cm}|}}
\hline
\rule[-0.2cm]{0pt}{0.6cm}\textbf{А} & \textbf{Б} & \textbf{В} & \textbf{Г} & \textbf{Д} \\
\hline
% Підпорка додана до кожного варіанту для ідеального вирівнювання
\rule[-0.4cm]{0pt}{0.9cm}#1 & 
\rule[-0.4cm]{0pt}{0.9cm}#2 & 
\rule[-0.4cm]{0pt}{0.9cm}#3 & 
\rule[-0.4cm]{0pt}{0.9cm}#4 & 
\rule[-0.4cm]{0pt}{0.9cm}#5 \\
\hline
\end{tabular}
}

% --- сумісність зі старою базою (діє, якщо тема не має власного визначення) ---
\providecommand{\trigStyle}[1]{\textbf{#1}}
\providecommand{\answerTableSmall}[5]{%
\begin{tabular}{|*{5}{>{\centering\arraybackslash}m{1.65cm}|}}
\hline
\rule[-0.2cm]{0pt}{0.6cm}\textbf{А} & \textbf{Б} & \textbf{В} & \textbf{Г} & \textbf{Д} \\
\hline
\rule[-0.4cm]{0pt}{0.9cm}#1 & \rule[-0.4cm]{0pt}{0.9cm}#2 & \rule[-0.4cm]{0pt}{0.9cm}#3 & \rule[-0.4cm]{0pt}{0.9cm}#4 & \rule[-0.4cm]{0pt}{0.9cm}#5 \\
\hline
\end{tabular}}
\providecommand{\matchingLayout}[3]{%
    \noindent
    \begin{minipage}[t]{0.40\textwidth}
        #1
    \end{minipage}%
    \hfill
    \begin{minipage}[t]{0.28\textwidth}
        #2
    \end{minipage}%
    \hfill
    \begin{minipage}[t]{0.30\textwidth}
        \vspace{0pt}
        \begin{flushright}
        #3
        \end{flushright}
    \end{minipage}}
\providecommand{\matchingLayoutBottom}[2]{%
    \noindent
    \begin{minipage}[t]{0.48\textwidth}
        #1
    \end{minipage}%
    \hfill
    \begin{minipage}[t]{0.48\textwidth}
        #2
    \end{minipage}
    \vspace{0.5cm}
    \noindent
    \matchingGrid}
\providecommand{\answerListVertical}[5]{%
    \vspace{0.2cm}
    \begin{itemize}[itemsep=0.4cm, leftmargin=1.5cm, labelsep=0.5cm]
        \item[\textbf{А}] #1
        \item[\textbf{Б}] #2
        \item[\textbf{В}] #3
        \item[\textbf{Г}] #4
        \item[\textbf{Д}] #5
    \end{itemize}
    \vspace{0.2cm}}

\chapterTitle{Тема 28a. Показникові рівняння}
\noindent{\small\color{gray!80} Розділ 2. Рівняння і нерівності \quad\textbullet\quad Усього завдань: 25 (НМТ 2023~--- 4, 2024~--- 7, 2025~--- 4, 2026~--- 10).}\par\vspace{0.2cm}

\typeTitle{НМТ 2023}

\noindent\zadnum Розв'яжіть систему рівнянь $\begin{cases} 2^{3x-y} = 2^7, \\ 2x + y = 3. \end{cases}$ Якщо $(x_0; y_0)$ --- розв'язок цієї системи, то $x_0 \cdot y_0 =$ \nmtyear{2023}

\answerTable{$2$}{$1$}{$-3$}{$-1$}{$-2$}

\noindent\zadnum Укажіть проміжок, якому належить корінь рівняння $\left(\displaystyle\frac{1}{3}\right)^{2x-1} = 9$. \nmtyear{2023}

\answerTable{$[2; +\infty)$}{$(-2; 0]$}{$(-\infty; -2]$}{$[1; 2)$}{$(0; 1]$}

\noindent\zadnum Розв'яжіть рівняння $3^x = 81 \cdot 9^x$. \nmtyear{2023}

\answerTable{$4$}{$-4$}{$-2$}{$-3$}{$3$}

% === ЗАВДАННЯ ===
\noindent\zadnum Якому проміжку належить корінь рівняння $2^x = \displaystyle\frac{1}{8}$? \nmtyear{2023}

\answerTable{$(2; 4]$}{$(-6; -4]$}{$(-2; 0]$}{$(0; 2]$}{$(-4; -2]$}

% === НМТ 2024 ===

\typeTitle{НМТ 2024}

\noindent\zadnum Укажіть проміжок, якому належить корінь рівняння $4^x \cdot 5^x = \displaystyle\frac{1}{400}$. \nmtyear{2024}

\answerTable{$[2; +\infty)$}{$[-0{,}5; 2)$}{$[-2; -0{,}5)$}{$[-10; -2)$}{$(-\infty; -10)$}

\noindent\zadnum Укажіть проміжок, якому належить корінь рівняння $5^{3x} \cdot 5^{-2x} = \displaystyle\frac{1}{5}$. \nmtyear{2024}

\answerTable{$(-\infty; -1]$}{$(-1; -0{,}5]$}{$(-0{,}5; 0]$}{$(0{,}5; +\infty)$}{$(0; 0{,}5]$}

\noindent\zadnum Укажіть корінь рівняння $2^x = 2 \cdot 16$. \nmtyear{2024}

\answerTable{$7$}{$3$}{$6$}{$4$}{$5$}

\noindent\zadnum Укажіть проміжок, якому належить корінь рівняння $\left(\displaystyle\frac{1}{3}\right)^{x-7} = \sqrt{3}$. \nmtyear{2024}

\answerTable{$(7; 10]$}{$(-4; 4)$}{$(-6; -4]$}{$(4; 7]$}{$[-7; -6]$}

\noindent\zadnum Укажіть проміжок, якому належить \textit{менший} корінь рівняння $3^{x^2} = 81$. \nmtyear{2024}

\answerTable{$[-2; 0)$}{$(2; +\infty)$}{$[-3; -2)$}{$[0; 2)$}{$(-\infty; -3)$}

\noindent\zadnum Розв'яжіть систему рівнянь $\begin{cases} 2^x + 2y = 0, \\ 2^x - 4y = 6. \end{cases}$ \nmtyear{2024}

\answerTable{$(-3; 3)$}{$(1; -1)$}{$(3; -3)$}{$(-1; 1)$}{$(2; -2)$}

\noindent\zadnum Розв'яжіть рівняння $10^x = 0{,}1$. \nmtyear{2024}

\answerTable{$1$}{$-1$}{$0$}{$0{,}01$}{$-9{,}9$}

% === НМТ 2025 ===

\typeTitle{НМТ 2025}

\noindent\zadnum Розв'яжіть систему рівнянь $\begin{cases} 2^{y+2x} = 64, \\ 2x - y = 12. \end{cases}$ Якщо $(x_0; y_0)$ --- розв'язок системи, то $x_0 + y_0 =$ \nmtyear{2025}

\answerTable{$4{,}5$}{$1{,}5$}{$7{,}5$}{$-3$}{$9$}

\noindent\zadnum Розв'яжіть рівняння $3^{x+1} = 81$. \nmtyear{2025}

\answerTable{$5$}{$3$}{$-3$}{$4$}{$2$}

\noindent\zadnum Укажіть проміжок, якому належить корінь рівняння $\left(\displaystyle\frac{3}{8}\right)^x = \displaystyle\frac{8}{3}$. \nmtyear{2025}

\answerTable{$(-2; 0)$}{$(2; 4)$}{$(-6; -4)$}{$(-4; -2)$}{$(0; 2)$}

\noindent\zadnum Укажіть проміжок, якому належить корінь рівняння $\left(\displaystyle\frac{1}{3}\right)^{x-7} = \sqrt{3}$. \nmtyear{2025}

\answerTable{$(-7; -6)$}{$(-6; 0)$}{$(0; 1)$}{$(6; 7)$}{$(1; 6)$}

\typeTitle{НМТ 2026}

% НМТ 2026, сесія 23.05, завдання №15 --- основна тема: Системи рівнянь
\zadtask{Розв'яжіть систему рівнянь $\begin{cases} 5^x = \sqrt{5}, \\ \dfrac{y + 1}{2x + 3} = \dfrac{1 - 4x}{2}. \end{cases}$ Якщо $(x_0;\, y_0)$~--- розв'язок системи, то $x_0 + y_0 = $ \nmtyear{2026}}
\answerTable{$-2{,}5$}{$-1{,}5$}{$2{,}5$}{$-3{,}5$}{$1{,}5$}
% Відповідь: А
\ifshowsolutions\solution{З першого рівняння $5^x=5^{1/2}$, тож $x_0=0{,}5$. Підставимо: $\dfrac{y+1}{2\cdot 0{,}5+3}=\dfrac{1-4\cdot 0{,}5}{2}$, тобто $\dfrac{y+1}{4}=-\dfrac{1}{2}$, звідки $y+1=-2$, $y_0=-3$. Тоді $x_0+y_0=0{,}5-3=-2{,}5$. Відповідь: \textbf{А}.}\fi

% НМТ 2026, сесія 1.06, завдання №1
\zadtask{Позначте число, що є коренем рівняння $3 \cdot 3^x = 27$. \nmtyear{2026}}
\answerTable{$2$}{$3$}{$4$}{$9$}{$27$}
% Відповідь: А
\ifshowsolutions\solution{$3\cdot 3^x=27 \Rightarrow 3^{x+1}=3^3 \Rightarrow x+1=3 \Rightarrow x=2$. Відповідь: \textbf{А}.}\fi

% НМТ 2026, сесія 3.06, завдання №11
\zadtask{Якому проміжку належить корінь рівняння $2^{x+1} - 2^x = \dfrac{1}{8}$? \nmtyear{2026}}
\answerTable{$(-\infty;\, -5)$}{$[-5;\, -3)$}{$[-3;\, 0)$}{$[0;\, 5)$}{$[5;\, +\infty)$}
% Відповідь: В
\ifshowsolutions\solution{$2^{x+1} - 2^x = 2^x(2-1) = 2^x = \dfrac18 = 2^{-3}$, звідки $x = -3 \in [-3;\, 0)$. Відповідь: \textbf{В}.}\fi

% НМТ 2026, сесія 5.06, завдання №13
\zadtask{Якому проміжку належить корінь рівняння $5^x \cdot 20^x = 1000$? \nmtyear{2026}}
\answerTable{$(-\infty;\, 0)$}{$[0;\, 1{,}5)$}{$[1{,}5;\, 2)$}{$[2;\, 4)$}{$[4;\, +\infty)$}
% Відповідь: В
\ifshowsolutions\solution{$5^x\cdot 20^x=(5\cdot 20)^x=100^x=10^{2x}$, а $1000=10^3$. Отже $10^{2x}=10^3\Rightarrow 2x=3\Rightarrow x=1{,}5$, що належить $[1{,}5;\,2)$. Відповідь: \textbf{В}.}\fi

% НМТ 2026, сесія 8.06, завдання №7
\zadtask{Якому проміжку належить корінь рівняння $2^{2x-5} = 8$? \nmtyear{2026}}
\answerTable{$(-\infty; -2)$}{$[0; 2)$}{$[-2; 0)$}{$[2; 4)$}{$[4; +\infty)$}

% НМТ 2026, сесія 9.06, завдання №7
\zadtask{Розв'яжіть рівняння $3^{x} = 81 \cdot 9^{x}$. \nmtyear{2026}}
\answerTable{$-2$}{$-4$}{$-3$}{$3$}{$4$}

% НМТ 2026, сесія 10.06, завдання №7
\zadtask{Позначте проміжок, якому належить корінь рівняння $\left(\dfrac{1}{3}\right)^{2x-1} = 9$. \nmtyear{2026}}
\answerTable{$(-3; -2)$}{$(-2; -1)$}{$(-1; 0)$}{$(0; 1)$}{$(1; 2)$}

% НМТ 2026, сесія 17.06, завдання №7
\zadtask{Розв'яжіть рівняння $3^x = 9^{x+2}$. \nmtyear{2026}}
\answerTable{$-4$}{$4$}{$0$}{$-2$}{$2$}
% Відповідь: А

% НМТ 2026, сесія 18.06, завдання №15 --- основна тема: Системи рівнянь
\zadtask{Розв'яжіть систему рівнянь $\begin{cases}2^{3x - y} = 2^7,\\ 2x + y = 3.\end{cases}$ Якщо $(x_0;\, y_0)$~--- розв'язок цієї системи, то $x_0 \cdot y_0 = $ \nmtyear{2026}}
\answerTable{$2$}{$-2$}{$-1$}{$-3$}{$1$}
% Відповідь: Б

% НМТ 2026, сесія 19.06, завдання №15
\zadtask{Визначте проміжок, якому належить корінь рівняння $5^{6x + 2} = \dfrac{1}{25}$. \nmtyear{2026}}
\answerTable{$[-2; -1)$}{$[-1; 0)$}{$[0; 1)$}{$[1; 2)$}{$[2; +\infty)$}
% Відповідь: Б

\ifshowsolutions
\sectionTitle{ВІДПОВІДІ ДО ЗАВДАНЬ НМТ 2026}
\begin{multicols}{3}\noindent
\noindent\textbf{16.}~А \ {\scriptsize\color{gray}(23.05, №15)}\par
\noindent\textbf{17.}~А \ {\scriptsize\color{gray}(1.06, №1)}\par
\noindent\textbf{18.}~В \ {\scriptsize\color{gray}(3.06, №11)}\par
\noindent\textbf{19.}~В \ {\scriptsize\color{gray}(5.06, №13)}\par
\noindent\textbf{23.}~А \ {\scriptsize\color{gray}(17.06, №7)}\par
\noindent\textbf{24.}~Б \ {\scriptsize\color{gray}(18.06, №15)}\par
\noindent\textbf{25.}~Б \ {\scriptsize\color{gray}(19.06, №15)}\par
\end{multicols}
\fi

\endgroup
\end{document}
